% --- Archon physics formalization source begin ---
% archon:physics
% archon:covers PhyXMiniProblems/problem_phyx_mini_0725.lean
% archon:source-report reports/phyx_mini/problem_phyx_mini_0725.source.json

\chapter{Physics problem phyx\_mini\_0725}
\label{ch:PhyXMiniProblems_problem_phyx_mini_0725}

\paragraph{Problem source.}
\#\# Physical scenario

Small hydroelectric plants in the mountains sometimes bring the water from a reservoir down to the power plant through enclosed tubes. In one such plant, the \textbackslash{}(100\textbackslash{},\textbackslash{}mathrm\{cm\}\textbackslash{})-diameter intake tube in the base of the dam is \textbackslash{}(50\textbackslash{},\textbackslash{}mathrm\{m\}\textbackslash{}) below the reservoir surface. The water drops \textbackslash{}(200\textbackslash{},\textbackslash{}mathrm\{m\}\textbackslash{}) through the tube before flowing into the turbine through a \textbackslash{}(50\textbackslash{},\textbackslash{}mathrm\{cm\}\textbackslash{})-diameter nozzle.

\#\# Question

By how much does the inlet pressure differ from the hydrostatic pressure at that depth?

\#\# Answer choices

- A: A: 1.1atm

- B: B: 1.3atm

- C: C: 1.5atm

- D: D: 1.7atm

\#\# Figure caption (auxiliary; use the image as primary evidence)

The image is a diagram illustrating a water flow system involving a dam and a turbine. It includes the following elements:

- **Dam**: A vertical structure is shown on the left side, with water stored on one side.
- **Water Levels**: The height of the water is marked at 250 meters, with a lower point inside the dam at 200 meters.
- **Points**: Three key points are labeled:
  - **Point 1**: Located at the top of the water level inside the dam.
  - **Point 2**: Positioned at the entrance to a streamline or pipe after the water drops 100 cm.
  - **Point 3**: At the end of the streamline, located where water exits or flows into the turbine, 50 cm above the turbine level.
  
- **Streamline**: A curved pathway or pipe indicating the flow of water from the dam to the turbine.

- **Turbine**: Positioned at the outflow end of the streamline, where the energy from the water flow is utilized.

- **Y-Axis (y in meters)**: Indicates the vertical scale in meters for height measurements.

The diagram effectively shows how water flows from the dam, through the streamline, and into the turbine, marking specific heights and changes in elevation along the way.

\#\# Dataset metadata

Dynamics; Physical Model Grounding Reasoning, Multi-Formula Reasoning

\paragraph{Recorded answer/context.}
C: C: 1.5atm

\paragraph{Figure/image path.}
/root/proposal\_for\_physic/science-mango/phyx\_mini\_run/phyx\_data/test\_image/725.png

\paragraph{Formalization target.}
create a compiling Lean file with sorry bodies at `PhyXMiniProblems/problem\_phyx\_mini\_0725.lean`. The Lean declarations must preserve the physical quantities, dimensions or dimensional roles, figure labels, governing-law hypotheses, and final relation expressed by this problem.
Use Mathlib/Physlib names found through LeanExplore where available. If a physics API is missing, introduce faithful local abstractions rather than scalar placeholder aliases.

\begin{theorem}[Physics formalization target]
\label{thm:physics:phyx_mini_0725:target}
\lean{PhyXMiniProblems.ProblemPhyXMini0725.problem_phyx_mini_0725}
\uses{def:physics:phyx-mini-0725:phyxminiproblems-problemphyxmini0725-hydroelectricsetup, def:physics:phyx-mini-0725:phyxminiproblems-problemphyxmini0725-inletpressuredeficitinpascals, def:physics:phyx-mini-0725:phyxminiproblems-problemphyxmini0725-matcheshydroelectricscenario, def:physics:phyx-mini-0725:phyxminiproblems-problemphyxmini0725-matchessupplieddamturbinefigure, def:physics:phyx-mini-0725:phyxminiproblems-problemphyxmini0725-matchesproblemdistancereadouts, def:physics:phyx-mini-0725:phyxminiproblems-problemphyxmini0725-usesstandardreferencedata, def:physics:phyx-mini-0725:phyxminiproblems-problemphyxmini0725-satisfiesreservoirandnozzleboundaryconditions, def:physics:phyx-mini-0725:phyxminiproblems-problemphyxmini0725-satisfiesidealhydroelectricflowlaws, def:physics:phyx-mini-0725:phyxminiproblems-problemphyxmini0725-answerchoice, def:physics:phyx-mini-0725:phyxminiproblems-problemphyxmini0725-displayedpressuredifferenceinatmospheres, def:physics:phyx-mini-0725:phyxminiproblems-problemphyxmini0725-roundstodisplayedtenthatmosphere, def:physics:phyx-mini-0725:phyxminiproblems-problemphyxmini0725-isnearestdisplayedpressuredifference}
The assigned autoformalize agent should translate this physics problem into Lean declarations in the covered file, with theorem and lemma proof bodies written as `by sorry`.
\end{theorem}
\begin{proof}
This is an autoformalization task, not a proof task. Produce statements that can later be proved by the physics prover without weakening the physical model.
\end{proof}

% --- Archon declaration topology begin ---
\section{Declaration topology}
\begin{definition}[Length Quantity]
  \label{def:physics:phyx-mini-0725:phyxminiproblems-problemphyxmini0725-lengthquantity}
  \lean{PhyXMiniProblems.ProblemPhyXMini0725.LengthQuantity}
  A nonnegative physical length, independent of a choice of units
\end{definition}
\begin{proof}
  This is the declared model definition of Length Quantity.
\end{proof}
\begin{definition}[Density Quantity]
  \label{def:physics:phyx-mini-0725:phyxminiproblems-problemphyxmini0725-densityquantity}
  \lean{PhyXMiniProblems.ProblemPhyXMini0725.DensityQuantity}
  A nonnegative mass density carrying dimension M L⁻³
\end{definition}
\begin{proof}
  This is the declared model definition of Density Quantity.
\end{proof}
\begin{definition}[Acceleration Quantity]
  \label{def:physics:phyx-mini-0725:phyxminiproblems-problemphyxmini0725-accelerationquantity}
  \lean{PhyXMiniProblems.ProblemPhyXMini0725.AccelerationQuantity}
  A nonnegative acceleration magnitude carrying dimension L T⁻²
\end{definition}
\begin{proof}
  This is the declared model definition of Acceleration Quantity.
\end{proof}
\begin{definition}[length In Meters]
  \label{def:physics:phyx-mini-0725:phyxminiproblems-problemphyxmini0725-lengthinmeters}
  \lean{PhyXMiniProblems.ProblemPhyXMini0725.lengthInMeters}
  \uses{def:physics:phyx-mini-0725:phyxminiproblems-problemphyxmini0725-lengthquantity}
  Read a physical length in coherent-SI metres
\end{definition}
\begin{proof}
  This is the declared model definition of length In Meters.
\end{proof}
\begin{definition}[density In Kilograms Per Cubic Meter]
  \label{def:physics:phyx-mini-0725:phyxminiproblems-problemphyxmini0725-densityinkilogramspercubicmeter}
  \lean{PhyXMiniProblems.ProblemPhyXMini0725.densityInKilogramsPerCubicMeter}
  \uses{def:physics:phyx-mini-0725:phyxminiproblems-problemphyxmini0725-densityquantity}
  Read a physical mass density in kilograms per cubic metre
\end{definition}
\begin{proof}
  This is the declared model definition of density In Kilograms Per Cubic Meter.
\end{proof}
\begin{definition}[acceleration In Meters Per Second Squared]
  \label{def:physics:phyx-mini-0725:phyxminiproblems-problemphyxmini0725-accelerationinmeterspersecondsquared}
  \lean{PhyXMiniProblems.ProblemPhyXMini0725.accelerationInMetersPerSecondSquared}
  \uses{def:physics:phyx-mini-0725:phyxminiproblems-problemphyxmini0725-accelerationquantity}
  Read a physical acceleration in metres per second squared
\end{definition}
\begin{proof}
  This is the declared model definition of acceleration In Meters Per Second Squared.
\end{proof}
\begin{definition}[pressure In Pascals]
  \label{def:physics:phyx-mini-0725:phyxminiproblems-problemphyxmini0725-pressureinpascals}
  \lean{PhyXMiniProblems.ProblemPhyXMini0725.pressureInPascals}
  Read a physical pressure in coherent-SI pascals
\end{definition}
\begin{proof}
  This is the declared model definition of pressure In Pascals.
\end{proof}
\begin{definition}[speed In Meters Per Second]
  \label{def:physics:phyx-mini-0725:phyxminiproblems-problemphyxmini0725-speedinmeterspersecond}
  \lean{PhyXMiniProblems.ProblemPhyXMini0725.speedInMetersPerSecond}
  Read a physical speed in coherent-SI metres per second
\end{definition}
\begin{proof}
  This is the declared model definition of speed In Meters Per Second.
\end{proof}
\begin{definition}[Flow Point]
  \label{def:physics:phyx-mini-0725:phyxminiproblems-problemphyxmini0725-flowpoint}
  \lean{PhyXMiniProblems.ProblemPhyXMini0725.FlowPoint}
  The three numbered points marked on the supplied dam-to-turbine figure
\end{definition}
\begin{proof}
  This is the declared model definition of Flow Point.
\end{proof}
\begin{definition}[Figure Object]
  \label{def:physics:phyx-mini-0725:phyxminiproblems-problemphyxmini0725-figureobject}
  \lean{PhyXMiniProblems.ProblemPhyXMini0725.FigureObject}
  Physical objects explicitly visible in the supplied raster
\end{definition}
\begin{proof}
  This is the declared model definition of Figure Object.
\end{proof}
\begin{definition}[Figure Label]
  \label{def:physics:phyx-mini-0725:phyxminiproblems-problemphyxmini0725-figurelabel}
  \lean{PhyXMiniProblems.ProblemPhyXMini0725.FigureLabel}
  Literal labels or readout annotations explicitly visible in the raster
\end{definition}
\begin{proof}
  This is the declared model definition of Figure Label.
\end{proof}
\begin{definition}[Dam Turbine Figure]
  \label{def:physics:phyx-mini-0725:phyxminiproblems-problemphyxmini0725-damturbinefigure}
  \lean{PhyXMiniProblems.ProblemPhyXMini0725.DamTurbineFigure}
  \uses{def:physics:phyx-mini-0725:phyxminiproblems-problemphyxmini0725-lengthquantity, def:physics:phyx-mini-0725:phyxminiproblems-problemphyxmini0725-flowpoint, def:physics:phyx-mini-0725:phyxminiproblems-problemphyxmini0725-figureobject, def:physics:phyx-mini-0725:phyxminiproblems-problemphyxmini0725-figurelabel}
  Geometry and qualitative content transcribed from the primary figure
\end{definition}
\begin{proof}
  This is the declared model definition of Dam Turbine Figure.
\end{proof}
\begin{definition}[Working Fluid]
  \label{def:physics:phyx-mini-0725:phyxminiproblems-problemphyxmini0725-workingfluid}
  \lean{PhyXMiniProblems.ProblemPhyXMini0725.WorkingFluid}
  The liquid named by the problem
\end{definition}
\begin{proof}
  This is the declared model definition of Working Fluid.
\end{proof}
\begin{definition}[Hydroelectric Setup]
  \label{def:physics:phyx-mini-0725:phyxminiproblems-problemphyxmini0725-hydroelectricsetup}
  \lean{PhyXMiniProblems.ProblemPhyXMini0725.HydroelectricSetup}
  \uses{def:physics:phyx-mini-0725:phyxminiproblems-problemphyxmini0725-lengthquantity, def:physics:phyx-mini-0725:phyxminiproblems-problemphyxmini0725-densityquantity, def:physics:phyx-mini-0725:phyxminiproblems-problemphyxmini0725-accelerationquantity, def:physics:phyx-mini-0725:phyxminiproblems-problemphyxmini0725-flowpoint, def:physics:phyx-mini-0725:phyxminiproblems-problemphyxmini0725-damturbinefigure, def:physics:phyx-mini-0725:phyxminiproblems-problemphyxmini0725-workingfluid}
  The independent quantities and idealizations of the hydroelectric setup. Neither hydrostaticPressureAtIntake nor the flowing inlet pressure is defined from an answer choice
\end{definition}
\begin{proof}
  This is the declared model definition of Hydroelectric Setup.
\end{proof}
\begin{definition}[bernoulli Energy Density In Pascals]
  \label{def:physics:phyx-mini-0725:phyxminiproblems-problemphyxmini0725-bernoullienergydensityinpascals}
  \lean{PhyXMiniProblems.ProblemPhyXMini0725.bernoulliEnergyDensityInPascals}
  \uses{def:physics:phyx-mini-0725:phyxminiproblems-problemphyxmini0725-lengthinmeters, def:physics:phyx-mini-0725:phyxminiproblems-problemphyxmini0725-densityinkilogramspercubicmeter, def:physics:phyx-mini-0725:phyxminiproblems-problemphyxmini0725-accelerationinmeterspersecondsquared, def:physics:phyx-mini-0725:phyxminiproblems-problemphyxmini0725-pressureinpascals, def:physics:phyx-mini-0725:phyxminiproblems-problemphyxmini0725-speedinmeterspersecond, def:physics:phyx-mini-0725:phyxminiproblems-problemphyxmini0725-flowpoint, def:physics:phyx-mini-0725:phyxminiproblems-problemphyxmini0725-hydroelectricsetup}
  Bernoulli energy density at one labelled point, read in pascals
\end{definition}
\begin{proof}
  This is the declared model definition of bernoulli Energy Density In Pascals.
\end{proof}
\begin{definition}[hydrostatic Reference Pressure In Pascals]
  \label{def:physics:phyx-mini-0725:phyxminiproblems-problemphyxmini0725-hydrostaticreferencepressureinpascals}
  \lean{PhyXMiniProblems.ProblemPhyXMini0725.hydrostaticReferencePressureInPascals}
  \uses{def:physics:phyx-mini-0725:phyxminiproblems-problemphyxmini0725-lengthinmeters, def:physics:phyx-mini-0725:phyxminiproblems-problemphyxmini0725-densityinkilogramspercubicmeter, def:physics:phyx-mini-0725:phyxminiproblems-problemphyxmini0725-accelerationinmeterspersecondsquared, def:physics:phyx-mini-0725:phyxminiproblems-problemphyxmini0725-pressureinpascals, def:physics:phyx-mini-0725:phyxminiproblems-problemphyxmini0725-hydroelectricsetup}
  The pressure predicted at the intake by a static water column measured from the open reservoir surface
\end{definition}
\begin{proof}
  This is the declared model definition of hydrostatic Reference Pressure In Pascals.
\end{proof}
\begin{definition}[inlet Pressure Deficit In Pascals]
  \label{def:physics:phyx-mini-0725:phyxminiproblems-problemphyxmini0725-inletpressuredeficitinpascals}
  \lean{PhyXMiniProblems.ProblemPhyXMini0725.inletPressureDeficitInPascals}
  \uses{def:physics:phyx-mini-0725:phyxminiproblems-problemphyxmini0725-pressureinpascals, def:physics:phyx-mini-0725:phyxminiproblems-problemphyxmini0725-hydroelectricsetup}
  The hydrostatic-minus-flowing inlet pressure difference, in pascals
\end{definition}
\begin{proof}
  This is the declared model definition of inlet Pressure Deficit In Pascals.
\end{proof}
\begin{definition}[inlet Pressure Deficit In Atmospheres]
  \label{def:physics:phyx-mini-0725:phyxminiproblems-problemphyxmini0725-inletpressuredeficitinatmospheres}
  \lean{PhyXMiniProblems.ProblemPhyXMini0725.inletPressureDeficitInAtmospheres}
  \uses{def:physics:phyx-mini-0725:phyxminiproblems-problemphyxmini0725-pressureinpascals, def:physics:phyx-mini-0725:phyxminiproblems-problemphyxmini0725-hydroelectricsetup, def:physics:phyx-mini-0725:phyxminiproblems-problemphyxmini0725-inletpressuredeficitinpascals}
  The same pressure deficit expressed in standard atmospheres
\end{definition}
\begin{proof}
  This is the declared model definition of inlet Pressure Deficit In Atmospheres.
\end{proof}
\begin{definition}[Matches Hydroelectric Scenario]
  \label{def:physics:phyx-mini-0725:phyxminiproblems-problemphyxmini0725-matcheshydroelectricscenario}
  \lean{PhyXMiniProblems.ProblemPhyXMini0725.MatchesHydroelectricScenario}
  \uses{def:physics:phyx-mini-0725:phyxminiproblems-problemphyxmini0725-hydroelectricsetup}
  Qualitative idealizations stated or conventionally used by the model
\end{definition}
\begin{proof}
  This is the declared model definition of Matches Hydroelectric Scenario.
\end{proof}
\begin{definition}[Matches Supplied Dam Turbine Figure]
  \label{def:physics:phyx-mini-0725:phyxminiproblems-problemphyxmini0725-matchessupplieddamturbinefigure}
  \lean{PhyXMiniProblems.ProblemPhyXMini0725.MatchesSuppliedDamTurbineFigure}
  \uses{def:physics:phyx-mini-0725:phyxminiproblems-problemphyxmini0725-lengthinmeters, def:physics:phyx-mini-0725:phyxminiproblems-problemphyxmini0725-flowpoint, def:physics:phyx-mini-0725:phyxminiproblems-problemphyxmini0725-figureobject, def:physics:phyx-mini-0725:phyxminiproblems-problemphyxmini0725-figurelabel, def:physics:phyx-mini-0725:phyxminiproblems-problemphyxmini0725-hydroelectricsetup}
  Primary-image evidence and the compatible dimensions stated in the prose. The figure places points 1, 2, and 3 at 250 m, 200 m, and 0 m, labels the intake and nozzle diameters as 100 cm and 50 cm, and shows one streamline leading from the reservoir through the dam to the turbine
\end{definition}
\begin{proof}
  This is the declared model definition of Matches Supplied Dam Turbine Figure.
\end{proof}
\begin{definition}[Matches Problem Distance Readouts]
  \label{def:physics:phyx-mini-0725:phyxminiproblems-problemphyxmini0725-matchesproblemdistancereadouts}
  \lean{PhyXMiniProblems.ProblemPhyXMini0725.MatchesProblemDistanceReadouts}
  \uses{def:physics:phyx-mini-0725:phyxminiproblems-problemphyxmini0725-lengthinmeters, def:physics:phyx-mini-0725:phyxminiproblems-problemphyxmini0725-hydroelectricsetup}
  The two vertical distances explicitly stated in the problem prose
\end{definition}
\begin{proof}
  This is the declared model definition of Matches Problem Distance Readouts.
\end{proof}
\begin{definition}[Uses Standard Reference Data]
  \label{def:physics:phyx-mini-0725:phyxminiproblems-problemphyxmini0725-usesstandardreferencedata}
  \lean{PhyXMiniProblems.ProblemPhyXMini0725.UsesStandardReferenceData}
  \uses{def:physics:phyx-mini-0725:phyxminiproblems-problemphyxmini0725-densityinkilogramspercubicmeter, def:physics:phyx-mini-0725:phyxminiproblems-problemphyxmini0725-accelerationinmeterspersecondsquared, def:physics:phyx-mini-0725:phyxminiproblems-problemphyxmini0725-pressureinpascals, def:physics:phyx-mini-0725:phyxminiproblems-problemphyxmini0725-hydroelectricsetup}
  Standard reference values used for the elementary numerical estimate: liquid-water density 1000 kg/m³, g = 9.8 m/s², and one standard atmosphere 101325 Pa
\end{definition}
\begin{proof}
  This is the declared model definition of Uses Standard Reference Data.
\end{proof}
\begin{definition}[Satisfies Reservoir And Nozzle Boundary Conditions]
  \label{def:physics:phyx-mini-0725:phyxminiproblems-problemphyxmini0725-satisfiesreservoirandnozzleboundaryconditions}
  \lean{PhyXMiniProblems.ProblemPhyXMini0725.SatisfiesReservoirAndNozzleBoundaryConditions}
  \uses{def:physics:phyx-mini-0725:phyxminiproblems-problemphyxmini0725-pressureinpascals, def:physics:phyx-mini-0725:phyxminiproblems-problemphyxmini0725-speedinmeterspersecond, def:physics:phyx-mini-0725:phyxminiproblems-problemphyxmini0725-hydroelectricsetup}
  Boundary conditions at the large free surface and the nozzle outlet
\end{definition}
\begin{proof}
  This is the declared model definition of Satisfies Reservoir And Nozzle Boundary Conditions.
\end{proof}
\begin{definition}[Satisfies Ideal Hydroelectric Flow Laws]
  \label{def:physics:phyx-mini-0725:phyxminiproblems-problemphyxmini0725-satisfiesidealhydroelectricflowlaws}
  \lean{PhyXMiniProblems.ProblemPhyXMini0725.SatisfiesIdealHydroelectricFlowLaws}
  \uses{def:physics:phyx-mini-0725:phyxminiproblems-problemphyxmini0725-lengthinmeters, def:physics:phyx-mini-0725:phyxminiproblems-problemphyxmini0725-pressureinpascals, def:physics:phyx-mini-0725:phyxminiproblems-problemphyxmini0725-speedinmeterspersecond, def:physics:phyx-mini-0725:phyxminiproblems-problemphyxmini0725-flowpoint, def:physics:phyx-mini-0725:phyxminiproblems-problemphyxmini0725-hydroelectricsetup, def:physics:phyx-mini-0725:phyxminiproblems-problemphyxmini0725-bernoullienergydensityinpascals, def:physics:phyx-mini-0725:phyxminiproblems-problemphyxmini0725-hydrostaticreferencepressureinpascals}
  Governing laws for steady, incompressible, inviscid flow: circular-tube continuity after cancellation of the common factor π/4; equality of Bernoulli energy density along the shown streamline; and the usual hydrostatic pressure law for the no-flow reference pressure. No numerical pressure deficit or answer choice occurs in these laws
\end{definition}
\begin{proof}
  This is the declared model definition of Satisfies Ideal Hydroelectric Flow Laws.
\end{proof}
\begin{definition}[Answer Choice]
  \label{def:physics:phyx-mini-0725:phyxminiproblems-problemphyxmini0725-answerchoice}
  \lean{PhyXMiniProblems.ProblemPhyXMini0725.AnswerChoice}
  Labels of the four pressure-difference choices printed in the problem
\end{definition}
\begin{proof}
  This is the declared model definition of Answer Choice.
\end{proof}
\begin{definition}[displayed Pressure Difference In Atmospheres]
  \label{def:physics:phyx-mini-0725:phyxminiproblems-problemphyxmini0725-displayedpressuredifferenceinatmospheres}
  \lean{PhyXMiniProblems.ProblemPhyXMini0725.displayedPressureDifferenceInAtmospheres}
  \uses{def:physics:phyx-mini-0725:phyxminiproblems-problemphyxmini0725-answerchoice}
  The four displayed pressure differences, expressed in atmospheres
\end{definition}
\begin{proof}
  This is the declared model definition of displayed Pressure Difference In Atmospheres.
\end{proof}
\begin{definition}[Rounds To Displayed Tenth Atmosphere]
  \label{def:physics:phyx-mini-0725:phyxminiproblems-problemphyxmini0725-roundstodisplayedtenthatmosphere}
  \lean{PhyXMiniProblems.ProblemPhyXMini0725.RoundsToDisplayedTenthAtmosphere}
  \uses{def:physics:phyx-mini-0725:phyxminiproblems-problemphyxmini0725-hydroelectricsetup, def:physics:phyx-mini-0725:phyxminiproblems-problemphyxmini0725-inletpressuredeficitinatmospheres}
  The calculated deficit rounds to a displayed tenth of an atmosphere
\end{definition}
\begin{proof}
  This is the declared model definition of Rounds To Displayed Tenth Atmosphere.
\end{proof}
\begin{definition}[Is Nearest Displayed Pressure Difference]
  \label{def:physics:phyx-mini-0725:phyxminiproblems-problemphyxmini0725-isnearestdisplayedpressuredifference}
  \lean{PhyXMiniProblems.ProblemPhyXMini0725.IsNearestDisplayedPressureDifference}
  \uses{def:physics:phyx-mini-0725:phyxminiproblems-problemphyxmini0725-hydroelectricsetup, def:physics:phyx-mini-0725:phyxminiproblems-problemphyxmini0725-inletpressuredeficitinatmospheres, def:physics:phyx-mini-0725:phyxminiproblems-problemphyxmini0725-answerchoice, def:physics:phyx-mini-0725:phyxminiproblems-problemphyxmini0725-displayedpressuredifferenceinatmospheres}
  A displayed choice is at least as close as every alternative
\end{definition}
\begin{proof}
  This is the declared model definition of Is Nearest Displayed Pressure Difference.
\end{proof}
% --- Archon declaration topology end ---
% --- Archon physics formalization source end ---
